% ============================================================
% SUPPLEMENTARY FILE for:
% "Integrating a FAIR-Compliant Early-Life PBK Model into FAIRDatabase:
%  Implementation and Web Interface for PFAS Exposure Assessment
%  in Children"
%
% Submitted to: Computational Toxicology
% Special Issue: Advancing FAIR and Integrative PBK Modelling for
%                Next-Generation Risk Assessment
% Article type in Editorial Manager: 'VSI: FAIR and PBK'
%
% NOTE FOR AUTHORS:
% This supplementary file contains all implementation-level detail that
% was moved out of the main manuscript to maintain readability and focus.
% The following sections are included:
%   S1.  SBML module descriptions (moved from main text Table 1)
%   S2.  Complete listing of 64 SBML state variables
%   S3.  Complete REST API endpoint table (moved from main text Table 3)
%   S4.  Chemical parameters for PFOA and PFOS
%   S5.  Full breastfeeding scenario parameterization
%   S6.  Installation and dependency instructions
%   S7.  Python runner code listing (simulate_scipy.py excerpt)
%   S8.  Role-based access control detailed table
%   S9.  User interface detailed specifications
%
% Upload this file as "supplementary_material.pdf" or "supplementary.tex"
% through Editorial Manager alongside the main manuscript.
% ============================================================

\documentclass[12pt]{article}
\usepackage{graphicx}
\usepackage{booktabs}
\usepackage{array}
\usepackage{multirow}
\usepackage{amsmath}
\usepackage{listings}
\usepackage{xcolor}
\usepackage{url}
\usepackage{enumitem}
\usepackage{hyperref}
\usepackage[margin=2.5cm]{geometry}

% Code listing style
\lstset{
  basicstyle=\ttfamily\footnotesize,
  keywordstyle=\color{blue}\bfseries,
  commentstyle=\color{gray},
  stringstyle=\color{orange},
  breaklines=true,
  frame=single,
  numbers=left,
  numberstyle=\tiny\color{gray}
}

\title{\textbf{Supplementary File}\\[0.5em]
\large Integrating a FAIR-Compliant Early-Life PBK Model into
FAIRDatabase: Implementation and Web Interface for PFAS Exposure
Assessment in Children}

\author{Tu-Ky Ly, Senna Wiedeman, Roland V. Bumbuc, Vivek M. Sheraton}
\date{}

\begin{document}
\maketitle
\tableofcontents
\newpage

% ============================================================
% S1. SBML MODULE DESCRIPTIONS
% ============================================================
\section*{S1. SBML Module Descriptions}
\addcontentsline{toc}{section}{S1. SBML Module Descriptions}

Table~\ref{tab:sbml_modules} describes the Python source modules that
compose the SBML encoding of the lifetime PBK model. Each module
generates a logically distinct component of the SBML file
\texttt{lifetime\_pbpk.xml}.

\begin{table}[h]
\centering
\caption{Python source modules composing the SBML encoding of the
lifetime PBPK model.}
\label{tab:sbml_modules}
\begin{tabular}{p{3.8cm} p{9.5cm}}
\toprule
\textbf{Module} & \textbf{Description} \\
\midrule
\texttt{compartments.py}
& Compartment definitions for all maternal and foetal organs. \\
\texttt{species.py}
& Declaration of the 64 state variables representing PFAS
  amounts (mg) per compartment. \\
\texttt{parameters.py}
& All constants and per-simulation inputs, including
  tissue-to-plasma partition coefficients (\texttt{PC\_*}),
  plasma protein binding parameters (albumin and
  $\alpha_1$-acid glycoprotein), plasma elimination half-life
  (\texttt{HalfLife}), and dietary intake rate
  (\texttt{RateInj}). \\
\texttt{rules\_time.py}
& Assignment rules for simulation time, subject age, and
  gestational timing. \\
\texttt{rules\_growth.py}
& A 5-phase, sex-dependent body weight model valid from birth
  to age 85, anchored at birth weight with $C^1$ continuity
  enforced at phase boundaries. \\
\texttt{rules\_foetus.py}
& Foetal weight trajectory, haematocrit, plasma protein
  binding, and foetal organ volumes and flows. \\
\texttt{rules\_volumes.py}
& Maternal organ volume rules. \\
\texttt{rules\_flows.py}
& Cardiac output, organ blood flows, milk production flow,
  and oral intake. \\
\texttt{rules\_chemistry.py}
& Tissue concentrations, enzymatic metabolism (liver, lung,
  kidney, gut, placenta), renal elimination rate
  (\texttt{Kelim}), and respiratory parameters. \\
\texttt{rules\_odes.py}
& The 64 rate rules (ODEs) governing PFAS mass balance in
  each compartment. \\
\bottomrule
\end{tabular}
\end{table}

PFAS toxicokinetics are governed by multi-compartment distribution
using plasma partition coefficients, with elimination via renal
excretion, biliary excretion, hepatic and pulmonary metabolism,
lactation, and a plasma-level first-order elimination pathway.
All quantities use consistent units: time in minutes, PFAS amounts
in mg, concentrations in mg/L, flows in L/min, body weights in kg.
Age in years is derived as $\text{Age} = t / 525{,}960$, where $t$
is solver time in minutes.

% ============================================================
% S2. COMPLETE LISTING OF SBML STATE VARIABLES
% ============================================================
\section*{S2. Complete Listing of SBML State Variables}
\addcontentsline{toc}{section}{S2. Complete Listing of SBML State Variables}

% NOTE FOR AUTHORS: Populate this table with the actual 64 species
% names from species.py. The structure below shows the grouping.

The SBML model comprises 64 state variables (species) organized
into two physiological sub-models:

\textbf{Maternal sub-model (23 compartments):} adipose, liver, kidney
(cortex, medulla), lung, gut, brain, breast, uterus, placenta,
amniotic fluid, venous blood pool, arterial blood pool, muscle, bone,
skin, heart, spleen, gonads, slowly perfused tissues.

\textbf{Foetal sub-model (19 compartments):} foetal liver, foetal
kidney, foetal lung, foetal brain, foetal gut, foetal blood, foetal
adipose, foetal muscle, foetal bone, foetal skin, foetal heart,
foetal spleen, foetal placenta-facing compartment, and auxiliary
mass-balance accumulators.

\textbf{Mass-balance accumulators (22 terms):} cumulative urinary
excretion, faecal excretion, lactation loss, hepatic metabolism,
oral ingestion, and auxiliary time-series terms.

% ============================================================
% S3. REST API ENDPOINT TABLE
% ============================================================
\section*{S3. REST API Endpoints}
\addcontentsline{toc}{section}{S3. REST API Endpoints}

Table~\ref{tab:api} lists all HTTP endpoints exposed by the
PBKFAIRModel Flask blueprint registered under the \texttt{/model}
URL prefix in FAIRDatabase.

\begin{table}[h]
\centering
\caption{API endpoints exposed by the PBKFAIR blueprint in FAIRDatabase.}
\label{tab:api}
\begin{tabular}{p{1.2cm} p{5.5cm} p{5.5cm} p{2cm}}
\toprule
\textbf{Method} & \textbf{Path} & \textbf{Description} &
\textbf{Role required} \\
\midrule
GET  & \texttt{/model/ui}
     & Simulation UI page
     & Login \\
POST & \texttt{/model/run}
     & Execute a scenario; returns JSON summary + time-series
     & Login \\
GET  & \texttt{/model/scenarios}
     & List available scenarios
     & Public \\
POST & \texttt{/model/parameter-sets}
     & Store a named parameter set
     & Admin/Curator \\
GET  & \texttt{/model/parameter-sets}
     & List all public parameter sets
     & Login \\
GET  & \texttt{/model/parameter-sets/\{id\}}
     & Fetch one parameter set with full params
     & Login \\
POST & \texttt{/model/runs}
     & Create and execute a simulation run
     & Admin/Curator \\
GET  & \texttt{/model/runs/\{id\}}
     & Retrieve a simulation run result
     & Accessor \\
POST & \texttt{/model/runs/\{id\}/artifacts}
     & Upload an artifact (image, CSV, mesh)
     & Admin/Curator \\
GET  & \texttt{/model/runs/\{id\}/artifacts}
     & List run artifacts
     & Accessor \\
DELETE & \texttt{/model/artifacts/\{id\}}
       & Delete an artifact
       & Admin/Owner \\
\bottomrule
\end{tabular}
\end{table}

A simulation is triggered via POST to \texttt{/model/run} with a JSON
body:
\begin{lstlisting}[language=Python, caption={Example POST request body
for /model/run}]
{
    "scenario": "bf_1yr",
    "HalfLife": 2.5,
    "RateInj": 0.451695,
    "BirthYear": 2007
}
\end{lstlisting}

% ============================================================
% S4. CHEMICAL PARAMETERS FOR PFOA AND PFOS
% ============================================================
\section*{S4. Chemical Parameters for PFOA and PFOS}
\addcontentsline{toc}{section}{S4. Chemical Parameters for PFOA and PFOS}

\begin{table}[h]
\centering
\caption{Reference chemical parameters for PFOA and PFOS in the
PBK model.}
\label{tab:chem_params}
\begin{tabular}{lcc}
\toprule
\textbf{Parameter} & \textbf{PFOA} & \textbf{PFOS} \\
\midrule
Plasma half-life (years)       & 2.5  & 3.5   \\
Free plasma fraction           & 0.02 & 0.006 \\
\texttt{HalfLife} (model var.) & 2.5  & 3.5   \\
Reference dietary intake (ng/kg/min) & 0.451695 & -- \\
\bottomrule
\end{tabular}
\end{table}

% NOTE FOR AUTHORS: Add full partition coefficients (PC_adipose,
% PC_liver, PC_kidney, etc.) for PFOA and PFOS from parameters.py.

% ============================================================
% S5. BREASTFEEDING SCENARIO PARAMETERIZATION
% ============================================================
\section*{S5. Breastfeeding Scenario Parameterization}
\addcontentsline{toc}{section}{S5. Breastfeeding Scenario Parameterization}

\begin{table}[h]
\centering
\caption{Detailed parameterization for the four reference breastfeeding
scenarios (male child, birth year 2007, PFOA).}
\label{tab:scenario_params}
\begin{tabular}{lcccc}
\toprule
\textbf{Parameter} & \textbf{no\_bf} & \textbf{bf\_6mo} &
\textbf{bf\_1yr} & \textbf{bf\_3yr} \\
\midrule
\texttt{StopBreastmilk\_total} (min)
  & 0 & 262,800 & 525,600 & 1,576,800 \\
\texttt{LactationTotal\_Duration\_inWeek}
  & 0 & 26 & 52 & 156 \\
\texttt{Gestation\_StartAge\_inYear}
  & -- & -- & -- & -- \\
\texttt{C\_milk\_input} (mg/L)
  & 0 & [value] & [value] & [value] \\
\bottomrule
\end{tabular}
\end{table}

% NOTE FOR AUTHORS: Fill in the C_milk_input and
% Gestation_StartAge_inYear values from the scenario definitions.

Pregnancy scenarios are activated by setting the mother's age at
conception (\texttt{Gestation\_StartAge\_inYear}), gestational
duration (\texttt{Gestation\_Duration\_inWeek} = 40 weeks), and
lactation duration (\texttt{LactationTotal\_Duration\_inWeek}).

% ============================================================
% S6. INSTALLATION AND DEPENDENCY INSTRUCTIONS
% ============================================================
\section*{S6. Installation and Dependency Instructions}
\addcontentsline{toc}{section}{S6. Installation and Dependency Instructions}

The PBKFAIRModel integration requires four additional Python packages
appended to \texttt{FAIRDatabase/backend/requirements.txt}:

\begin{lstlisting}[language=bash, caption={Additional Python dependencies}]
python-libsbml>=5.19
scipy>=1.11
numpy
pandas
\end{lstlisting}

The Flask blueprint is registered in \texttt{app.py} by adding:

\begin{lstlisting}[language=Python, caption={Blueprint registration in
FAIRDatabase app.py}]
from src.model.routes import model_routes
app.register_blueprint(model_routes, url_prefix="/model")
\end{lstlisting}

For production deployment, it is strongly recommended to wrap the
\texttt{run\_scenario()} call in an asynchronous background worker
(e.g., Celery with Redis) and poll for results, as synchronous
execution blocks the HTTP request thread for 30--90 seconds per
scenario.

The complete package structure is:

\begin{lstlisting}[language=bash, caption={PBKFAIRModel package structure
within the FAIRDatabase repository}]
FAIRDatabase/
  PBKFAIRModel/           # Self-contained Python package
    lifetime_pbpk.xml     # Committed SBML model (source of truth)
    runner.py             # PBPKModel class with execute() API
    sbml/                 # SBML generation modules
      generate_sbml.py
      compartments.py
      species.py
      parameters.py
      rules_time.py
      rules_growth.py
      rules_foetus.py
      rules_volumes.py
      rules_flows.py
      rules_chemistry.py
      rules_odes.py
  backend/src/model/      # Flask blueprint
    routes.py             # API endpoints
    helpers.py            # Parameter validation and dispatch
  frontend/templates/model/
    pbkfair.html          # Jinja2 simulation UI template
\end{lstlisting}

% ============================================================
% S7. PYTHON RUNNER CODE LISTING
% ============================================================
\section*{S7. Python Runner Code Listing}
\addcontentsline{toc}{section}{S7. Python Runner Code Listing}

The Python simulation runner exposes the following public API:

\begin{lstlisting}[language=Python, caption={PBKFAIRModel public API
(simulate\_scipy.py)}]
from PBKFAIRModel import execute

result = execute({
    "scenario": "bf_1yr",   # One of: no_bf, bf_6mo, bf_1yr, bf_3yr
    "HalfLife": 2.5,         # Plasma half-life in years (PFOA default)
    "RateInj": 0.451695,     # Dietary intake rate in ng/kg/min
    "BirthYear": 2007        # Subject birth year
})

# Result structure:
# result["peak_C_ven"]   -- Peak venous plasma concentration (mg/L)
# result["age_at_peak"]  -- Age at peak (years)
# result["final_C_ven"]  -- Venous concentration at end (mg/L)
# result["final_age"]    -- Age at end of simulation (years)
# result["timeseries"]   -- DataFrame, 500 downsampled time points
#   Columns: time (min), Age (yr), BDW (kg), C_ven (mg/L),
#            C_art (mg/L), C_foetus_22 (mg/L), C_milk (mg/L),
#            + mass-balance accumulators
\end{lstlisting}

% ============================================================
% S8. ROLE-BASED ACCESS CONTROL
% ============================================================
\section*{S8. Role-Based Access Control}
\addcontentsline{toc}{section}{S8. Role-Based Access Control}

\begin{table}[h]
\centering
\caption{Role-based access control (RBAC) for PBKFAIRModel endpoints
in FAIRDatabase.}
\label{tab:rbac}
\begin{tabular}{p{2.5cm} p{3.5cm} p{3.5cm} p{3.5cm}}
\toprule
\textbf{Permission} & \textbf{Admin} & \textbf{Curator} &
\textbf{Accessor} \\
\midrule
View simulation UI          & \checkmark & \checkmark & \checkmark \\
Run scenarios               & \checkmark & \checkmark & \checkmark \\
Create parameter sets       & \checkmark & \checkmark & $\times$ \\
Create simulation runs      & \checkmark & \checkmark & $\times$ \\
Upload artifacts            & \checkmark & \checkmark & $\times$ \\
View parameter sets         & \checkmark & \checkmark & \checkmark \\
View simulation runs        & \checkmark & \checkmark & \checkmark \\
Delete any artifact         & \checkmark & $\times$   & $\times$ \\
Delete own artifact         & \checkmark & \checkmark & $\times$ \\
\bottomrule
\end{tabular}
\end{table}

% ============================================================
% S9. USER INTERFACE DETAILED SPECIFICATIONS
% ============================================================
\section*{S9. User Interface Detailed Specifications}
\addcontentsline{toc}{section}{S9. User Interface Detailed Specifications}

\subsection*{S9.1 Layout and Navigation}

The simulation interface extends the FAIRDatabase dashboard layout
(Jinja2 template: \texttt{frontend/templates/model/pbkfair.html}),
maintaining visual consistency with the broader platform. The page is
accessible at \texttt{/model/ui} and requires user authentication.
Four main sections: parameter configuration, simulation execution,
results visualization, and artifact management.

\subsection*{S9.2 Visualization}

Results are rendered using Chart.js as a line chart of venous PFAS
concentration (mg/L) versus subject age (years). Interactive features
include hover tooltips, zoom and pan. A colour-blind-friendly palette
(blue for concentration) is used with clear axis labels and units.

\subsection*{S9.3 Artifact Storage}

Supported upload formats: JPG, PNG, MP4, VTK, VTU, VTP (max 200 MB).
Storage backend: Supabase Storage with signed URLs (10-minute TTL)
for secure temporary access. Features include a real-time upload
progress bar, inline image and video preview, one-click download with
original filename preservation, and role-gated deletion.

\subsection*{S9.4 CSV Export}

A download button exports the complete time-series as a CSV file
including all output variables (time, age, concentrations, body
weights, mass-balance terms) at 500-point downsampled resolution.

\end{document}
